\section[缺陷与缺陷算子]{缺陷与缺陷算子（4 月 20 日）}

\srcpage{65}

设 $Z$ 是 $n$ 维拓扑场论。算子-态对应表明，局部算子等价于 $Z(S^{n-1})$ 中的态；在算子乘积展开下，只要 $n\ge2$，这个向量空间就是交换代数。此外，对任何闭 $(n-1)$-流形 $M$，$Z(M)$ 自然地是 $Z(S^{n-1})$-模。而对闭 $(n-2)$-流形 $N$，$Z(N)$ 不只是 $\C$-线性范畴，它还自然地是 $Z(S^{n-1})$-线性的：挖去一个小球，就得到 $Z$ 的边界 $S^{n-1}$ 作用在 Hom 空间上的方式。

然后令 $\cM^0_Z$（即我们此前称为 $\cM^{\mathrm{aff}}_Z$ 的东西）为（仿射）vacua 模空间，即 $\Spec Z(S^{n-1})$。$Z$ 赋给流形的东西层化到 $\cM^0_Z$ 上。我们想把 $Z$ 与一个映入 $\cM^0_Z$ 的映射的 Lagrange 场论联系起来。例如对规范群为 $G$ 的二维有限规范理论，$\cM^0_Z=\hat G$，这让我们把 $Z$ 分解为不可约表示。

为了更抽象地做到这一点，有一个在某种意义上最容易写下的理论会有帮助。设 $R$ 是交换环，$X=\Spec R$。于是有一个任意维的拓扑场论，称为 B-模型（B-model），以这些数据为基础。不过它是截断的：我们不给最高维流形赋数，只给低维对象赋不变量。想法是：这是一个 Lagrange 场论，场是映射 $M\to X$。这有点怪——$M$ 是光滑流形而 $X$ 是仿射簇，所以我们做一个天真的处理，要求这些映射局部常值。于是它们等价于映射 $\pi_0(M)\to X$。这是后面更有趣版本的这个理论的玩具模型。换种说法，$\operatorname{Map}(\pi_0(M),X)=\Spec R\otimes\pi_0(M)$。这不依赖于维数或别的什么，这要么很好（灵活！），要么很无聊。

现在必须线性化，在选定的维数 $n$ 定义截断 TFT $B^n_X$。对闭 $(n-1)$-流形 $M$ 与闭 $(n-2)$-流形 $N$，分别赋
\begin{align*}
M&\longmapsto \cO(\operatorname{Map}(\pi_0(M),X))=R\otimes\pi_0(M),\\
N&\longmapsto \QCoh(\operatorname{Map}(\pi_0(N),X))=\mathrm{Mod}_{R\otimes\pi_0(N)}.
\end{align*}
可以定义协边上的推拉运算。

现在设 $Z$ 是维数 $2$ 或以上的任意 TFT，考虑其 vacua 模空间上的 B-模型 $B^n_{\cM^0_Z}$。我们想猜 $Z$ 被这个 B-模型很好地逼近。一般来说这远非事实，但对二维有限规范理论成立：$Z_G(S^1)=\C[G/G]=\C[\hat G]$，这正是 B-模型所赋的东西；例如 $Z_G(\mathrm{pt})=\Rep G$ 等价于 B-模型所赋的 $\QCoh(\hat G)$。

这是一个非常天真的构造，但用更精细的方式做，它导向镜像对称（mirror symmetry）！这最接近 Gross--Siebert 的内蕴镜像对称。从一个有趣的场论（A-模型）出发，把簇（镜像）构造为局部算子的谱，然后镜像对称把 A-模型与 B-模型等同起来。细节不平凡：必须用导出的方式工作，有时还要修改镜像，但这是可以做到的。不过即使现在也有很多要修的：上次我们看到，对由范畴 $\Rep G(R)$ 定义的二维 TFT，这行不通。

取而代之，我们得到的是：对 $M$ 中任一点，$Z(M)$ 上有 $Z(S^{n-1})$-模结构，所以我们知道该在哪里挖球。两点之间的路径给出它们的模结构之间的同构；但要点是这个条件只是常值的，所以 $Z(M)$ 自然地只是 $\pi_0(M)\otimes Z(S^{n-1})$ 上的模，即 $\cO(\operatorname{Map}(\pi_0(M),\cM^0_Z))$ 上的模。也就是说，$Z(M)$ 是这个映射空间上的层，即 B-模型高一维时赋给 $M$ 的对象：$Z(M)\in B^{n+1}_{\cM^0_Z}(M)$。类似地，赋给 $(n-2)$-流形的不变量，对 $N$ 中任选一点都是 $Z(S^{n-1})$-线性范畴；一般地这只是局部常值的，所以又要与 $\pi_0(N)$ 作张量。特别地，我们得到 2-范畴 $B_{\cM^0_Z}(N)$ 中的一个对象，即 $\operatorname{Map}(\pi_0(N),\cM^0_Z)$ 上的范畴层，或者说是 $\pi_0(N)\otimes Z(S^{n-1})$-线性的范畴。

某种意义下，这是一般的态-可观测量对应（state-observable correspondence）：态是 $Z$ 的输出，可观测量是高维 B-模型的输出。简言之，$Z$ 典范地是 $B^{n+1}_{\cM^0_Z}$ 的一个边界条件。这在结构上组织了局部算子及其谱信息给出的全部数据。

在四维，我们研究了 $G=\mathrm{GL}_1$ 的「A-模型」，即拓扑 Maxwell 理论。我们说明过 $\cA_{\mathrm{GL}_1}(Y)$（$Y$ 是三维流形）是 $Y$ 上 $U_1$-联络模规范等价的空间上的局部常值函数空间。我们想理解这个理论中的局部算子。为此回忆场空间分解为 $H^2(Y;\Z)$ 与一个环面的乘积。对局部算子，我们看

\srcpage{66}

$Y=S^3$，此时 $H^1$ 与 $H^2$ 都消失——所以没有有趣的东西。再仔细一点，你实际上有 $BU_1$，但 $S^2\to BU_1$ 上没有任何有趣的局部常值函数。

所以看来出路是过渡到导出世界，在那里得到 $H^*(BU_1)=\C[u]$，这不是世上最有趣的代数。不幸——但当我们把局部算子推广到更有趣的可观测量（缺陷）时，事情会变得更有趣。

这就把我们带到缺陷与缺陷算子，我们将从协边假设（cobordism hypothesis）的角度接近它们。这是一个把拓扑场论中的概念与范畴论中的概念联系起来的定理。可以用它从范畴论数据造 TFT；或者反过来，用它把「用图像推理高阶范畴理论」的方式精确化。协边假设编码了用于推理各类范畴的所有不同的图形演算，并解释了它们为何存在。

设 $\Box$ 是一个切结构（tangential structure），意思是我们可以放在流形切丛上的一种拓扑结构——无结构、定向、spin 结构、标架等等。可以把它精确定义为群同态 $H_n\to O_n$，以及标架丛的结构群到主 $O_n$-丛的约化。\yizhu{底本在这里用方框 $\Box$ 作为切结构的一般记号，并非排版缺字；$\Box$-结构在协边范畴中充当标架、定向等的统一替代。}{A（原文 p.66）}有 $\Box$-结构的协边概念，于是得到对称幺半协边 $n$-范畴 $\cB\mathrm{ord}^{\Box}_n$。在维数小于 $n$ 的流形上，我们要求 $n$ 维领邻域上有 $\Box$-结构。我们把 $\Box$-流形的 TFT 定义为对称幺半函子
\[
Z:\cB\mathrm{ord}^{\Box}_n\longrightarrow\cC,
\]
其中 $\cC$ 是另一个对称幺半 $n$-范畴。

协边假设说：这样的 TFT 完全由对象 $Z(\mathrm{pt})\in\cC$ 决定（其中 $\mathrm{pt}$ 带典范 $\Box$-结构），但须满足
\begin{enumerate}
\item 有限性条件（可对偶性，dualizability），以及
\item 某种自对偶结构，由 $\Box$ 指定。
\end{enumerate}
有限性条件可以想成发芽（germination）：从你希望赋给一个点的东西出发，把场论向上生长；你需要很强的有限性假设才能一路长上去。

协边假设最初由 Baez--Dolan [BD95] 猜想，然后由 Hopkins--Lurie 在一维（TODO：一维还是二维？）以及 Lurie [Lur09] 在全维数证明。\footnote{还有 Schommer--Pries 与 Ayala--Francis 的部分证明。}

Baez--Dolan [BD95] 把它推广为另一个猜想，称为缠结假设（tangle hypothesis）；Lurie [Lur09, §4.3] 进一步推广并证明了所谓带奇点的协边假设（cobordism hypothesis with singularities）。其中我们描述带各种指定奇点的流形上的场论。例如可以固定流形 $M$ 并允许 $M$ 上的锥；于是我们的协边范畴看起来像带奇点的 $\Box$-流形，奇点的链环（link）是 $M$（TODO：我可能理解错了）。

于是带奇点的协边假设说：从修改后的协边范畴出发的对称幺半函子，等价于此前那样的 TFT $Z$，连同 $Z(M)$ 的一个对象；更确切地说一个映射 $1\to Z(M)$，不管 $M$ 维数多少都有意义。这就是对称幺半函子赋给 $M$ 的锥的东西，而锥是从 $\mathrm{pt}$ 到 $M$ 的协边。对此也有有限性与自对偶条件。

我们最感兴趣的是 $M=S^k$ 的情形。例如 $k=n-1$ 时，带奇点的流形是带标记点的 $n$-流形，而带奇点的协边假设给这份数据赋一个通常的 TFT $Z$，连同 $Z(S^{n-1})$ 的一个元素。也就是说，我们挑出了一个可以放在标记点处的局部算子。更一般地，若我们关心带嵌入的 $(k-1)$ 维子流形的流形，链环是 $S^{n-k}$，于是带奇点的协边假设说：给定 $Z$ 与 $Z(S^{n-k})$ 中的一个对象，就可以计算这类流形。这就是我们称为缺陷或算子的东西。例如线缺陷或线算子由范畴 $Z(S^{n-2})$ 的对象标记。也就是说，$Z(S^{n-2})$ 是线算子的范畴——按定义这是一个范畴，但也可以从拓扑上看出：两个线算子之间的态射就是它们之间的界面，即可以放在红线段与蓝线段交接处的东西。

\srcpage{67}

同样地，$Z(S^{n-3})$ 是面缺陷的 2-范畴：有面缺陷以及它们之间的态射（放在一维交接处），而这些还可以有态射（放在交接处之间的零维交接处）。轮中有轮……

现在这些嵌入子流形可以打结。例如若 $Z$ 是三维 TFT 并选一个线算子，我们可以在带嵌入纽结的三维流形上计算 $Z$，得到一个数。这件事（及其到高维缺陷的推广）就是带奇点的协边假设允许我们做的。

工具箱里的算子更多了，来看看我们还有什么额外结构。我们用局部算子代数定义 $\cM^0_Z$，并对理论中的对象作谱分解。对线算子，同样有算子乘积展开的概念，它把 $Z(S^{n-2})$ 变成幺半范畴。在 TFT 中，这来自挖去两个小球的球所给出的协边
\[
S^{n-2}\amalg S^{n-2}\longrightarrow S^{n-2}.
\]
我们把它想成发生在线缺陷上的一个点处，即线在此点的链环。同样，小球嵌入大球的方式有不同选择，所以我们实际上在构形空间上以局部常值的方式得到这份数据。也就是说，$Z(S^{n-2})$ 有许多乘法，由更大圆盘内一对圆盘的构形空间标记。形式上，这是小 $(n-1)$-圆盘 operad 上的代数结构，即 $E_{n-1}$-结构；不过在用导出视角之前，这不怎么有用。

例如 $n=2$ 时，线算子是自界面，而算子乘积展开结合但不交换：前式是更大区间内一对区间的协边，这份数据的构形空间不连通。因此乘积 $L_1*L_2$ 与 $L_2*L_1$ 不能由构形空间中的道路连接，一般不等；但 $L_1*(L_2*L_3)$ 与 $(L_1*L_2)*L_3$ 可以连接，所以这个代数结合。一般地我们得到（范畴中的）$E_1$-代数。

若 $n=3$，乘法在更大圆盘内小圆盘的空间中是局部常值的。更精确地说，这意味着 TODO。构形空间有非平凡的 $\pi_1$，所以取一个非平凡环路 $\ell$。它定义了从 $L_1*L_2$ 到 $L_2*L_1$ 的映射，但我们得不到更多信息，例如不保证它平方等于恒等。这给 $Z(S^{n-2})$ 一个编织幺半范畴（braided monoidal category）的结构，或者说范畴中的 $E_2$-代数。

对 $n\ge4$，构形空间单连通，因此有自然同构 $L_1*L_2\to L_2*L_1$——这是 $E_3$-范畴，它是对称幺半的（事实上是 $E_\infty$-范畴）。于是在四维 TFT（以及我们不关心的更高维）中，可以定义 $\cM^1_Z$ 为带算子乘积展开的对称幺半范畴 $Z(S^2)$ 的 $\Spec$。这是 $\cM^0_Z$ 的细化，后者是它的仿射化（affinization）。高维有更好的几何，从而有更好的谱分解；这份额外结构对 Langlands 纲领的 TFT 进路很重要。这就是我们下次要做的，例如会看到四维 Maxwell 理论中有大量线算子。
